\section{Related Work}
\label{sec:related-work}

GAWorld v2 draws on five bodies of work that answer different parts of the
artificial-society problem.  Its contribution is a system integration and an
audited set of research surfaces; we do not claim that the underlying ideas
of memory, tool-using agents, social evaluation, agent-based validation, or
agent platforms are unique to GAWorld.

\subsection{Generative agents and simulated individuals}

Generative Agents demonstrated that an LLM-driven architecture combining an
experience stream, retrieval, reflection, and planning can sustain socially
legible behavior in a sandbox community \citep{park2023generative}.  Later
work grounds agents in extensive self-report data and evaluates how well
responses predict held-out answers from the same individuals
\citep{park2024grounded}.  Research on ``silicon samples'' explores
conditioning language models to approximate distributions of survey responses
\citep{argyle2023outofone}.  These studies motivate profile grounding and
longitudinal state, but they also distinguish two validation targets:
internally coherent simulated inhabitants and externally evaluated models of
particular people or populations.

GAWorld v2 is oriented toward constructing the surrounding research system.
It adds configurable urban, social, economic, work, organizational,
information, and policy mechanisms; multiple operational interfaces including
a class-scale multi-tenant deployment; and explicit experiment and artifact
workflows around a persistent-agent loop.  Its archived profiles and traces
have not undergone the held-out, person-level validation used in
self-report-grounded agents.  Accordingly, v2 adopts persistence as a system
property and treats fidelity as an open empirical claim, as the v1 system
paper already argued \citep{gawoerldv1}.

\subsection{Language-agent cognition and architecture}

LLM-agent research has developed reusable patterns for coupling language
generation to action.  ReAct interleaves reasoning and environment-facing
actions \citep{yao2023react}; Reflexion stores verbal feedback to improve
later attempts \citep{shinn2023reflexion}; Voyager combines iterative
prompting, environment feedback, and an accumulating skill library in an
open-ended embodied setting \citep{wang2023voyager}.  MemoryBank studies
long-term memory with retrieval and forgetting mechanisms
\citep{zhong2024memorybank}, while CoALA offers a conceptual organization of
language agents in terms of memory, action space, and decision procedures
\citep{sumers2024coala}.  CAMEL examines role-conditioned communication
among language agents \citep{li2023camel}.

v2 informs its separation of structured state, episodic storage, retrieval,
planning, action, reflection, and learned skills from these systems.  It
differs in scope rather than by proposing a new generic reasoning algorithm:
it connects such components to a persistent multi-agent society, a clocked
file-backed runtime, an event bus, a plugin registry, a per-tenant world
model, and a parallel-worlds engine.  v2 additionally adopts a
\emph{pluggable personality} layer so the cognitive profile of a resident can
be replaced without rewriting the rest of the loop.  The contribution is an
implemented integration for social simulation, not a claim of a universally
superior cognitive architecture.

\subsection{Social-agent tasks and evaluation}

AgentBench evaluates LLM agents across interactive environments and emphasizes
the gap between language ability and reliable action
\citep{liu2024agentbench}.  AgentBoard adds fine-grained progress measures
and analytical visualization for multi-turn agents \citep{ma2024agentboard}.
SOTOPIA focuses specifically on social intelligence by placing agents in
interactive social scenarios and evaluating goal achievement and social
dimensions \citep{zhou2024sotopia}.  More recent work proposes social
benchmarks that go beyond pairwise interaction, including GroupBench for
collaborative social reasoning and the social-cognition probes discussed in
\citet{hashimoto2022generate} and \citet{leng2024llm}.

v2 addresses a complementary unit of analysis: a persistent society run
whose agents, environment, institutions, and artifacts evolve across
simulated time.  Its benchmark layer can check selected outputs, known-sign
diagnostics, placebos, and determinism when suitable artifacts exist.  v2
adds two new evaluation surfaces: the \emph{Agent Playground} of fourteen
game-based probes, in which a small group of residents faces a structured
scenario under deterministic scoring rules; and the \emph{Agent Arena}, in
which a fixed set of residents answers a controlled question bank under
correctness and latency metrics.  Neither is a replacement for
standardized social-intelligence benchmarks; they are probes that
complement population-scale traces.

\subsection{Agent-based social simulation and validation}

Agent-based modeling (ABM) treats macro patterns as outcomes of specified
micro-level rules and interactions, while generative social science
emphasizes using such models to explain how collective regularities can
arise \citep{epstein2006generative}.  The ODD protocol provides a standard
vocabulary for documenting an individual- or agent-based model's overview,
design concepts, and details \citep{grimm2006odd}.  Methodological work on
empirical ABM validation stresses that matching selected aggregate
observations is not sufficient: calibration, identification, sensitivity,
multiple validation targets, and out-of-sample evidence remain distinct
concerns \citep{windrum2007empirical,fagiolo2007critical}.  The recent
\emph{Human-AI interaction} literature emphasizes the gap between model
behavior and population behavior when the same model is reused in different
contexts \citep{kapania2025between}.

v2 inherits this demand for explicit mechanisms and layered evidence, but
LLM mediation introduces additional sources of variation.  Prompt context,
provider behavior, model versions, parsing fallbacks, and generated text
can change trajectories even when conventional parameters and seeds are
recorded.  v2 therefore couples ABM-style configuration and intervention
design to provider logs, raw agent traces, per-arm completion status, an
artifact ledger, and an explicit evidence-class taxonomy.  Its present
archive supports software-path and affordance claims, not empirical
validation of a population model.  This distinction also responds to
concerns that language models should not be treated as drop-in replacements
for human participants \citep{dillion2023replace}, particularly when
generated outputs can flatten variation within identity groups
\citep{wang2025flatten}.

\subsection{Platforms, parallel execution, and counterfactual tooling}

A number of recent platforms have approached parts of the artificial-society
problem at production scale.  Concordia \citep{vezhnevets2023generative}
generalizes Generative Agents to a composable event-driven formulation with
a focus on embodied conversational characters.  OASIS \citep{yang2024oasis}
scales social-media-style simulation to one million agents.  AgentSociety
\citep{gao2025agentsociety} adds a city-scale multi-agent environment with
urban, economic, and information substrates.  RecSim \citep{ie2019recsim}
and its successors approach recommendation and counterfactual evaluation
from a different angle, treating each user as a model and each policy as an
intervention.  The research-engineering literature on parallel and
distributed simulation \citep{fujimoto2016parallel} emphasizes logical
processes, conservative synchronization, and rollback.  v2 does not adopt
conservative synchronization: the hub-node protocol is a best-effort relay
that allows local autonomy and tolerates brief disconnection, and the
parallel-worlds engine is sequential at the world level, not at the agent
level.  v2 instead positions itself as a platform whose goals are
configurability, auditability, and a single-cohort comparability across
matched worlds.

The counterfactual-inference layer (paired ATE with bootstrap CI, sign-flip
$p$, BH-$q$, pre-event balance / difference-in-differences, placebo noise
floor, heterogeneity, dose--response, replication across seeds, and an
optional one-call model reading) follows the parallel-worlds literature
\citep{pearl2018book,hernan2020causal,rubin2005causal} and the recent
multiple-testing literature \citep{benjamini1995fdr}.  It is deterministic
and does not call a model during inference; the model call is reserved for
the optional interpretation, which is itself bound to specific evidence rows.
